\documentclass{article}
\usepackage{arxiv}

\usepackage[utf8]{inputenc}
\usepackage[T1]{fontenc}
\usepackage{hyperref}
\usepackage{url}
\usepackage{booktabs}
\usepackage{amsmath, amssymb}
\usepackage{bm}
\usepackage{natbib}
\usepackage{graphicx}

% The emblem lives in the repository-wide assets directory, two levels up.
\graphicspath{{../../assets/}{../../assets/logos/}{../../assets/biota/}{../../assets/diagrams/}}

\title{The Mathematics and Design of DuskCrayfish}

% arxiv.sty labels every paper "A Preprint" in the running head and under the
% title. This is an internal technical report, not a preprint of a submission.
\renewcommand{\headeright}{Technical Report}
\renewcommand{\undertitle}{Technical Report}

% arxiv.sty's title block sets a 24pt strut above the author names and 0.4in
% below them, which leaves the byline stranded between "Technical Report" and
% the date. This is the same layout with those two gaps halved. Overridden here
% rather than edited in arxiv.sty, which is a vendored file.
\makeatletter
\renewcommand{\@maketitle}{%
  \vbox{%
    \hsize\textwidth
    \linewidth\hsize
    \vskip 0.1in
    \@toptitlebar
    \centering
    {\LARGE\sc \@title\par}
    \@bottomtitlebar
    \textsc{\undertitle}\\
    \vskip 0.1in
    \def\And{%
      \end{tabular}\hfil\linebreak[0]\hfil%
      \begin{tabular}[t]{c}\bf\rule{\z@}{12\p@}\ignorespaces%
    }
    \def\AND{%
      \end{tabular}\hfil\linebreak[4]\hfil%
      \begin{tabular}[t]{c}\bf\rule{\z@}{12\p@}\ignorespaces%
    }
    \begin{tabular}[t]{c}\bf\rule{\z@}{12\p@}\@author\end{tabular}%
  \vskip 0.2in \@minus 0.1in \center{\@date}   \vskip 0.2in
  }
}
\makeatother

\author{
  Daniel Boupha Christensen \quad Anisa Khair \quad Ian Liu \\
  Environmental, Health, and Safety, University of California, Berkeley \\
  The Green Initiative Fund \\
}

\begin{document}

% Masthead. Licence at the left, group emblem at the right. This is an internal
% report, so the block carries the licence and the project mark only, and
% deliberately no journal, volume, page range or DOI furniture.
\noindent
\begin{minipage}[c]{0.74\textwidth}
  \footnotesize
  Copyright 2024--2026 Strawberry Creek Monitoring Group contributors.\\
  Licensed under the Apache License, Version 2.0.\\
  \url{https://www.apache.org/licenses/LICENSE-2.0}
\end{minipage}%
\hfill
\begin{minipage}[c]{0.16\textwidth}
  \raggedleft
  \includegraphics[width=\linewidth]{SCMGemblem.png}
\end{minipage}

\maketitle

\begin{abstract}
DuskCrayfish learns the normal behaviour of a small urban creek and flags
moments it cannot predict well. Given a window of recent sensor history across
all the connected monitoring sites, it predicts the next reading at every site,
and a sharp divergence between the real and predicted conductivity is the
anomaly signal. We never show the model what a discharge looks like. It learns
only normal behaviour, and an anomaly is whatever it fails to reproduce. Here,
we set out the mathematics of the deployed model, a graph convolution over the
creek's flow topology followed by a per-node recurrent forecaster, together
with the detection rules we built on top of it and the reasoning behind each
choice. We kept the architecture small, and we give our reasons for this
alongside the mathematics. Finally, we describe how the model is served in the
StrawberryWatch dashboard and how the served checkpoint scored the labelled
events in October 2026.
\end{abstract}

\section{Introduction}
\label{sec:introduction}

The design of DuskCrayfish follows from one idea. We take a window of recent
sensor history across all the connected sites and predict the next reading at
every site. When the real reading and the prediction diverge sharply on
conductivity, that gap is the anomaly signal. The model is never told what a
discharge looks like; it learns only normal behaviour, and an anomaly is
whatever it fails to reproduce.

In this document we explain the mathematics and, alongside each piece, why we
built the model as we did. The architecture is small by design, and we give
the reasons at each point where the choice arises.

\subsection{The deployed configuration}
\label{sec:configuration}

Table~\ref{tab:config} lists the settings the model runs with, drawn from the
project's configuration. We refer to these values throughout the mathematics
that follows, so that the general equations stay tied to this specific model.

\begin{table}[t]
  \centering
  \caption{The deployed configuration. These are the settings the current
    checkpoint runs with, not defaults or recommendations.}
  \label{tab:config}
  \begin{tabular}{ll}
    \toprule
    Setting & Value \\
    \midrule
    Sequence length $T$ & 24 timesteps (6 hours at 15 min sampling) \\
    Features per node $F$ & 10 (3 sensor, 3 weather, 4 cyclical time) \\
    Hidden dimension $d$ & 16 \\
    Graph convolution layers & 1 \\
    LSTM layers & 1 \\
    Graph convolution type & GCN \\
    Dropout & 0.1 \\
    Training epochs & 10 (early stopping, patience 5) \\
    Batch size & 128 \\
    Learning rate & 0.001 \\
    Train / validation split & 0.8 / 0.2 \\
    Threshold percentile & 99 \\
    Rain lookback window & 12 hours \\
    Rain threshold multiplier $m$ & 2.0 \\
    Rain amount that counts as wet $\theta$ & 0.1 \\
    Post rain decay window & 36 hours \\
    Imputation gap limit & 24 hours \\
    \bottomrule
  \end{tabular}
\end{table}

One graph layer, one long short-term memory (LSTM) layer and a hidden
dimension of only 16 make a lean model, by design. The training data is a
rolling window of recent creek behavior rather than a large historical
archive, and a larger model trained on that much data would overfit,
memorizing the recent past rather than learning the creek's general rhythm. A
small model also trains in seconds on a CPU. This matters because the data is a
moving window, and we retrain the model regularly to keep it current rather
than training it once and freezing it. We sized the model to the data that
exists, not to what a larger model would suggest on paper.

\section{Data and Notation}
\label{sec:data}

We represent the creek as a graph with $N$ sensor nodes whose connections follow
the direction of water flow. At each timestep every node carries a vector of
features: (1) conductivity, depth and temperature from the sensors, (2) air
temperature, rainfall and solar radiation merged in from weather data when
available, and (3) four cyclical time encodings, the sine and cosine of hour of
day and of day of year. The number of features is not hardcoded. The model
reads the feature count from the data it receives at runtime and sizes its
input and output layers to match, which lets it adapt when a weather channel is
missing. We write $F$ for that count and treat $F = 10$ as the deployed case
throughout. Table~\ref{tab:notation} collects the symbols used in the rest of
this document.

\begin{table}[t]
  \centering
  \caption{Notation.}
  \label{tab:notation}
  \begin{tabular}{lll}
    \toprule
    Symbol & Meaning & Dimension \\
    \midrule
    $N$ & number of sensor nodes & scalar, $N = 4$ \\
    $F$ & features per node & scalar, $F = 10$ \\
    $T$ & timesteps in one input window & scalar, $T = 24$ \\
    $d$ & hidden dimension & scalar, $d = 16$ \\
    $d_e$ & node embedding dimension & scalar, $d_e = 8$ \\
    $\mathbf{X}$ & one input window & $\mathbb{R}^{T \times N \times F}$ \\
    $\mathbf{X}_t$ & feature snapshot at timestep $t$ & $\mathbb{R}^{N \times F}$ \\
    $\mathbf{Y}$ & target snapshot, the step after the window & $\mathbb{R}^{N \times F}$ \\
    $\hat{\mathbf{Y}}$ & predicted snapshot & $\mathbb{R}^{N \times F}$ \\
    $\mathbf{A}$ & directed flow adjacency & $\{0,1\}^{N \times N}$ \\
    $\hat{\mathbf{A}}$ & adjacency with self loops, $\mathbf{A} + \mathbf{I}$ & $\{0,1\}^{N \times N}$ \\
    $\hat{\mathbf{D}}$ & diagonal degree matrix of $\hat{\mathbf{A}}$ & $\mathbb{R}^{N \times N}$ \\
    $\mathbf{H}$ & node features at one timestep & $\mathbb{R}^{N \times F}$ \\
    $\mathbf{H}'$ & node features after graph convolution & $\mathbb{R}^{N \times d}$ \\
    $\mathbf{W}$ & graph convolution weights & $\mathbb{R}^{F \times d}$ \\
    $\mathbf{h}_t$ & pooled creek summary at timestep $t$ & $\mathbb{R}^{d}$ \\
    $\mathbf{e}_i$ & learned embedding of node $i$ & $\mathbb{R}^{d_e}$ \\
    $\mathbf{z}_{i,t}$ & temporal model input for node $i$ & $\mathbb{R}^{2d + d_e}$ \\
    $\mathbf{h}^{\text{lstm}}_{i,T}$ & final recurrent state for node $i$ & $\mathbb{R}^{d}$ \\
    $\mathbf{W}_{\text{out}}, \mathbf{b}_{\text{out}}$ & output layer & $\mathbb{R}^{F \times d}$, $\mathbb{R}^{F}$ \\
    $e_{ij}$ & prediction error, node $i$ feature $j$ & scalar \\
    $s_i$ & normalized anomaly score at node $i$ & scalar \\
    $\tau_i$ & alerting threshold at node $i$ & scalar \\
    $m, \theta, D$ & rain multiplier, wet threshold, decay window & scalars \\
    $k$ & level shift threshold in robust deviations & scalar, $k = 4.0$ \\
    \bottomrule
  \end{tabular}
\end{table}

A single training example is a window of $T = 24$ consecutive timesteps,

\begin{equation}
  \label{eq:window}
  \mathbf{X} \in \mathbb{R}^{T \times N \times F},
\end{equation}

where $\mathbf{X}_t \in \mathbb{R}^{N \times F}$ is the snapshot of every node's
features at timestep $t$. The target is the snapshot immediately after the
window:

\begin{equation}
  \label{eq:target}
  \mathbf{Y} = \mathbf{X}_{T+1} \in \mathbb{R}^{N \times F}.
\end{equation}

We chose $T = 24$ because it gives six hours of history at the creek's
fifteen minute cadence. This is long enough to capture the daily rhythms the
creek has (e.g., warming through the afternoon, overnight runoff from
sprinklers) without reaching so far back that old conditions outweigh the
recent trend the model is trying to extend.

Before anything else, we normalize every feature to a z-score using statistics
from complete data only:

\begin{equation}
  \label{eq:zscore}
  z = \frac{x - \mu}{\sigma}.
\end{equation}

The model relies on one consequence of this choice. After normalization, a
value of zero is the average, the neutral middle of the distribution. We use
the same value as the placeholder for missing sensor readings, so a gap
contributes nothing rather than pulling the model up or down. This lets the
model run on a network where some sites are offline: a sensor that is down
reads as the neutral zero wherever it is missing.

\subsection{The deployed four-node core}
\label{sec:graph}

We now describe the graph for the network the model runs on. The four active
nodes, in the order the code indexes them, are north fork 0, south fork 2,
south fork 1 and oxford. Directed edges connect them along the physical
direction of water flow downhill, and there are three of them:

\begin{equation}
  \label{eq:edges}
  \text{north fork 0} \rightarrow \text{oxford}, \quad
  \text{south fork 1} \rightarrow \text{south fork 2}, \quad
  \text{south fork 2} \rightarrow \text{oxford}.
\end{equation}

Writing these as a directed adjacency $\mathbf{A}$, where entry $(i,j)$ is one
when there is an edge from node $i$ to node $j$, in the code's index order (north
fork 0, south fork 2, south fork 1, oxford) gives

\begin{equation}
  \label{eq:adjacency}
  \mathbf{A} =
  \begin{pmatrix}
    0 & 0 & 0 & 1 \\
    0 & 0 & 0 & 1 \\
    0 & 1 & 0 & 0 \\
    0 & 0 & 0 & 0
  \end{pmatrix}.
\end{equation}

The rows of Equation~\ref{eq:adjacency} read as follows: north fork 0 flows to
oxford, south fork 2 flows to oxford, south fork 1 flows to south fork 2, and
oxford flows nowhere, since it is the most downstream point. The empty last row
marks a sink, the confluence where the network ends. The graph convolution then
adds a self loop to every node and applies the degree normalization described
in Section~\ref{sec:gcn}, so each node's update blends its own value with the
values flowing into it along these edges. The structure is sparse and
determined entirely by the creek's physical layout, so the model does not have
to learn which sites connect to which. We supply the connectivity directly, and
it is correct by construction.

\subsection{Exclusion of footbridge}
\label{sec:footbridge}

An earlier version of this network had five nodes, with footbridge between
north fork 0 and oxford on the north branch. We removed footbridge and
contracted its edge, so north fork 0 now flows directly to oxford. The water
still runs that way; there is no sensor in between.

We made this change because of a mismatch between training and deployment.
Inside the corpus window, footbridge has 1,191 real readings against 29,123
timesteps, a real coverage of 4.1 percent, and it does not appear in the live
data cache at all. All of its other values were imputed. A node that is almost
entirely imputed in training and entirely absent in production teaches the
model to rely on a signal that will not be there at run time, and this node
passed its imputed signal to its graph neighbours as though it were real.
Contracting the edge keeps the flow topology accurate and trains the model on
the same set of sites it serves.

\section{Model Architecture}
\label{sec:architecture}

\subsection{The spatial step, a graph convolution}
\label{sec:gcn}

At each timestep we let each node's features be informed by the nodes connected
to it, because a disturbance upstream appears downstream as the water carries
it. With $\mathbf{H} \in \mathbb{R}^{N \times F}$ the node features at one
timestep, the graph convolution \citep{kipf2017gcn} computes:

\begin{equation}
  \label{eq:gcn}
  \mathbf{H}' = \operatorname{ReLU}\!\left( \hat{\mathbf{D}}^{-\frac{1}{2}} \, \hat{\mathbf{A}} \, \hat{\mathbf{D}}^{-\frac{1}{2}} \, \mathbf{H} \, \mathbf{W} \right),
\end{equation}

where $\mathbf{W} \in \mathbb{R}^{F \times d}$ is a learned weight matrix
projecting the $F$ features into the hidden size $d = 16$, and the nonlinearity
is

\begin{equation}
  \label{eq:relu}
  \operatorname{ReLU}(x) = \max(0, x).
\end{equation}

This is the standard graph convolutional network (GCN) operator. The adjacency
with inserted self loops is $\hat{\mathbf{A}} = \mathbf{A} + \mathbf{I}$, and
$\hat{\mathbf{D}}$ is its diagonal degree matrix with $\hat{D}_{ii} = \sum_j
\hat{A}_{ij}$. The self loops let each node keep its own features rather than
seeing only its neighbours, and the degree scaling $\hat{\mathbf{D}}^{-1/2}$ on
both sides keeps a heavily connected site from dominating a lightly connected
one. The graph library \citep{fey2019pyg} adds the self loops and applies this
normalization automatically, so the edge list we give the model does not need
to include them.

The matrix form in Equation~\ref{eq:gcn} is compact, but it does not show what
happens at one site. Written per node, the new representation of node $i$ is a
normalized sum over itself and the nodes connected to it:

\begin{equation}
  \label{eq:gcn-per-node}
  \mathbf{h}'_i = \operatorname{ReLU}\!\left( \sum_{j \,\in\, \mathcal{N}(i) \,\cup\, \{i\}} \frac{1}{\sqrt{\hat{d}_i \, \hat{d}_j}} \, \mathbf{h}_j \, \mathbf{W} \right),
\end{equation}

where $\mathcal{N}(i)$ is the set of nodes connected to $i$ through the graph and
$\hat{d}_i$ is node $i$'s degree including its self loop. Thus, each node's
updated features are a degree-weighted blend of its own features and its
neighbours', passed through the learned projection $\mathbf{W}$ and the
nonlinearity. A spike in conductivity at one site therefore enters the
connected site's representation at the next layer, which is the propagation
the model needs to recognize a real event as it moves through the creek.

This step involves two design choices. The first is the type of convolution.
We set the GCN type explicitly in the configuration, and a GCN weights
neighbours by the fixed graph structure, which is the appropriate prior when
the connections are real water flow rather than relationships the model should
infer. One alternative would be to let the model learn its own neighbour
weightings. However, the flow topology is known physics, so there is little
reason to make the model learn weightings it would only overfit on limited
data. The second choice is depth. We use a single graph layer, so each node
sees its immediate neighbours. On a network this small, with short flow paths,
one hop already reaches the relevant sites, and stacking more layers would
blur every node toward a single averaged value, a known failure of deep graph
networks called oversmoothing \citep{li2018deeper}.

\subsection{Pooling into one creek state}
\label{sec:pooling}

The temporal model needs a single vector summarising the whole creek at each
timestep alongside the per-node detail, so we average across the nodes:

\begin{equation}
  \label{eq:pooling}
  \mathbf{h}_t = \frac{1}{N} \sum_{i=1}^{N} \mathbf{H}'_i \;\in\; \mathbb{R}^{d}.
\end{equation}

Doing this at every timestep produces a sequence $\mathbf{h}_1, \dots,
\mathbf{h}_T$, each entry a summary of the whole connected network at one
moment. We use mean pooling rather than a more elaborate combination because it
is simple and robust, and it works well with the missing-data handling: when a
node mask is supplied, we average over the reporting sites only, so the creek
summary reflects real data rather than placeholder zeros.

The pooled summary accompanies the per-node representations rather than
replacing them. An earlier version pooled and discarded $\mathbf{H}'$, which
left the temporal model unable to tell one site from another, so every node
received the same prediction. We carry both forward to the next step.

\subsection{The temporal step, an LSTM}
\label{sec:lstm}

The sequence enters a single-layer LSTM network \citep{hochreiter1997lstm},
which tracks how the signal evolves and carries relevant information forward
across steps. At each timestep it maintains a hidden state
$\mathbf{h}_t^{\text{lstm}}$ and a cell state $\mathbf{c}_t$, updated through
gates:

\begin{equation}
  \label{eq:lstm}
  \begin{aligned}
    \mathbf{f}_t &= \sigma\!\left( \mathbf{W}_f \mathbf{h}_t + \mathbf{U}_f \mathbf{h}_{t-1}^{\text{lstm}} + \mathbf{b}_f \right), \\
    \mathbf{i}_t &= \sigma\!\left( \mathbf{W}_i \mathbf{h}_t + \mathbf{U}_i \mathbf{h}_{t-1}^{\text{lstm}} + \mathbf{b}_i \right), \\
    \mathbf{o}_t &= \sigma\!\left( \mathbf{W}_o \mathbf{h}_t + \mathbf{U}_o \mathbf{h}_{t-1}^{\text{lstm}} + \mathbf{b}_o \right), \\
    \tilde{\mathbf{c}}_t &= \tanh\!\left( \mathbf{W}_c \mathbf{h}_t + \mathbf{U}_c \mathbf{h}_{t-1}^{\text{lstm}} + \mathbf{b}_c \right), \\
    \mathbf{c}_t &= \mathbf{f}_t \odot \mathbf{c}_{t-1} + \mathbf{i}_t \odot \tilde{\mathbf{c}}_t, \\
    \mathbf{h}_t^{\text{lstm}} &= \mathbf{o}_t \odot \tanh(\mathbf{c}_t).
  \end{aligned}
\end{equation}

The sigmoid $\sigma$ maps each of the three gates to values between zero and
one, so the gates act as soft switches. The symbol $\odot$ is elementwise
multiplication, and the $\mathbf{W}$, $\mathbf{U}$ and $\mathbf{b}$ are
learned. The forget gate $\mathbf{f}_t$ sets how much of the previous memory
to keep, the input gate $\mathbf{i}_t$ how much of the new candidate
$\tilde{\mathbf{c}}_t$ to store, and the output gate $\mathbf{o}_t$ how much of
the resulting memory to emit. The cell update line is where the network
retains or discards information.

We chose the LSTM over a plain recurrent network because its gates let it hold
a trend from hours earlier when it is relevant and discard it when it is not,
which matters for a six hour window. We use a single LSTM layer for the same
reason the rest of the model is small: the sequences are short and the data is
limited, so one layer captures the temporal pattern without overfitting. We
keep the final hidden state,

\begin{equation}
  \label{eq:final-state}
  \mathbf{h}_T^{\text{lstm}} \in \mathbb{R}^{d},
\end{equation}

which is the model's compressed summary of how the creek has behaved over the
whole window.

\subsection{The output, a prediction of the next reading}
\label{sec:output}

The temporal model runs per node rather than once for the whole creek. For node
$i$ at timestep $t$, the input is the concatenation of that node's own graph
convolution output, the pooled creek summary and a learned embedding of the
node's identity:

\begin{equation}
  \label{eq:lstm-input}
  \mathbf{z}_{i,t} = \left[\, \mathbf{H}'_{i,t} \;\Vert\; \mathbf{h}_t \;\Vert\; \mathbf{e}_i \,\right] \in \mathbb{R}^{2d + d_e},
\end{equation}

with $\mathbf{e}_i \in \mathbb{R}^{d_e}$, $d_e = 8$, one learned vector per
site. The node embedding lets the model represent, for example, that south fork
2 normally runs warmer than north fork 0, rather than treating every site as
interchangeable. The sequence $\mathbf{z}_{i,1}, \dots, \mathbf{z}_{i,T}$
passes through the LSTM of Equation~\ref{eq:lstm}, and the final hidden state
$\mathbf{h}_{i,T}^{\text{lstm}} \in \mathbb{R}^{d}$ is that node's own
compressed history.

The output layer predicts a change from the last observed reading rather than an
absolute value:

\begin{equation}
  \label{eq:output}
  \hat{\mathbf{Y}}_i = \mathbf{X}_{T,i} + \left( \mathbf{W}_{\text{out}} \, \mathbf{h}_{i,T}^{\text{lstm}} + \mathbf{b}_{\text{out}} \right),
\end{equation}

with $\mathbf{W}_{\text{out}} \in \mathbb{R}^{F \times d}$ and
$\mathbf{b}_{\text{out}} \in \mathbb{R}^{F}$ learned, and $\mathbf{X}_{T,i}$ the
node's features at the last timestep of the window, feature-propagated if a mask
filled it in. Everything is in normalized units, so no denormalization happens
at this step.

The residual form in Equation~\ref{eq:output} is what gives the model a notion
of where the readings currently are. One alternative would be to predict from
the hidden state alone. However, the model then falls back on whatever average
shape minimizes error across the training set, which is climatology rather
than a forecast.

Both choices are visible in the predictions. On the deployed checkpoint, the
spread across nodes, the largest node prediction minus the smallest, reaches
4.592 in normalized units. Of this spread, the learned head contributes 0.188
and the residual anchor contributes 4.591. Thus, almost all of the per-node
difference comes from each node starting at its own last reading, and the head
supplies a comparatively small correction on top. A version without node
identity would give every node the same number, a spread of exactly zero.

\section{Training}
\label{sec:training}

We train by minimizing the mean squared error between prediction and truth,

\begin{equation}
  \label{eq:loss}
  \mathcal{L} = \frac{1}{\lvert \mathcal{M} \rvert \, \lvert \mathcal{S} \rvert} \sum_{(i) \in \mathcal{M}} \; \sum_{j \in \mathcal{S}} \left( \hat{\mathbf{Y}}_{ij} - \mathbf{Y}_{ij} \right)^2,
\end{equation}

over training data that contains only ordinary creek behavior. The model thus
learns the absence of discharges, not discharges themselves.

Two restrictions narrow that sum. First, $\mathcal{S}$ is the scored channel
set (conductivity, depth and temperature) rather than all $F$ channels. The
model still predicts every channel it receives. However, rainfall 15 minutes
ahead cannot be learned from creek sensors, and scoring it adds irreducible
error that is largest during storms, when conductivity matters most. Second,
$\mathcal{M}$ is the set of node targets that were observed and passed quality
control. A target that was never read is filled with the neutral zero, and
scoring the model against that placeholder zero teaches it to reproduce an
absence.

Two further training choices are worth noting. We apply dropout of 0.1 during
training, zeroing a small random fraction of hidden units at each step, which
prevents the small model from relying too heavily on any single pattern.
Training runs for up to 10 epochs with early stopping at patience 5: if the
validation loss does not improve for five epochs, the run stops and the best
weights are restored, so the model does not train past the point where it is
still learning something general. We train in plain 32 bit precision on CPU,
because the mixed precision that speeds up large models on a graphics card
breaks the LSTM computation on a CPU and offers no benefit at this size.

\section{Detection}
\label{sec:detection}

\subsection{The model all, alert one strategy}
\label{sec:score}

Once trained, the model runs on new data, and we examine where it missed. The
prediction error at node $i$ on feature $j$ is the absolute deviation:

\begin{equation}
  \label{eq:error}
  e_{ij} = \left\lvert \hat{\mathbf{Y}}_{ij} - \mathbf{Y}_{ij} \right\rvert.
\end{equation}

The defining choice is which of these errors we act on. The model predicts all
features, but we use only the conductivity error to decide whether something
is anomalous, which we call the model all, alert one strategy. Conductivity is
the channel that responds to discharges, while the others, rainfall in
particular, produce large prediction errors unrelated to water quality. An
earlier version that scored on the average error across all features was in
effect flagging rainstorms, because the rain prediction error outweighed
everything else. Isolating conductivity removed those false alarms.

We keep the score per node rather than averaging it across the network. Writing
$c$ for the conductivity index, node $i$'s score at a timestep is $e_{ic}$, and
we judge it against that node's own history rather than a shared bar:

\begin{equation}
  \label{eq:score}
  s_i = \frac{e_{ic} - \operatorname{median}_i}{\operatorname{IQR}_i},
\end{equation}

where $\operatorname{median}_i$ and $\operatorname{IQR}_i$ are the median and
interquartile range (IQR) of node $i$'s conductivity error over held-out
validation data. The earlier design averaged across nodes, on the reasoning
that an event propagating through the network should score higher than a blip
at a single site. We replaced it because sites differ. A site that runs
noisier than the rest sits permanently above a shared threshold, and averaging
lets a large deviation at one node be diluted by three quiet ones. Robust
normalization per node addresses both problems, and keeping the node axis lets
an alert name the affected site.

The threshold is also per node, fixed once during training as the 99th
percentile of that node's normalized validation errors:

\begin{equation}
  \label{eq:threshold}
  \tau_i = \operatorname{percentile}_{99}\big( \{ s_i^{\text{val}} \} \big).
\end{equation}

We fix the threshold at training time rather than recomputing it on whatever
data is in front of the model. Recomputing would be circular and would let the
threshold drift with each batch. The 99th percentile means that under normal
conditions only about one percent of timesteps at a node would cross it, which
sets the false alarm rate before deployment rather than discovering it
afterwards. A single timestep over the bar is not an alert. A node is flagged
only when at least three scored timesteps in the window cross it, so one noisy
reading cannot raise an alarm.

We compute the median, IQR and threshold on observed targets only. A node
target that was never read is filled with the neutral zero, and including such
targets in the calibration sets the bar against data that does not exist. This
matters in practice. On the current corpus, 27.7 percent of one site's target
cells were imputed, and its threshold was calibrated almost entirely against
the error of predicting a zero.

\subsection{The second rule, sustained level shift}
\label{sec:rule2}

The error-based rule above cannot detect one shape. A clean step in
conductivity surprises the model once and is then absorbed into its own
forecast, so the residual goes quiet while the water stays abnormal. A second
rule therefore watches the level itself rather than the error:

\begin{equation}
  \label{eq:level-shift}
  u_i = \frac{\mathbf{Y}_{ic} - \operatorname{median}^{\text{lvl}}_i}{\operatorname{IQR}^{\text{lvl}}_i}, \qquad \lvert u_i \rvert > k,
\end{equation}

with $k = 4.0$, and with the median and IQR taken over the node's normalized
conductivity level on validation data. A node is anomalous if either rule
fires, and the result records which one did, since an onset and a sustained
abnormal level mean different things operationally.

\subsection{The rain adjustment}
\label{sec:rain}

Rain raises conductivity without being a discharge, as runoff carries road salt
and grime into the creek. To avoid flagging ordinary rain, we raise the
threshold during wet periods. Summing rainfall over the preceding 12 hour
window, the per-timestep threshold is

\begin{equation}
  \label{eq:rain}
  \tau_t =
  \begin{cases}
    m \cdot \tau & \text{if } \sum_{k \in \text{window}} r_k > \theta, \\[4pt]
    \tau \left( 1 + (m - 1)\left(1 - \operatorname{clip}\!\left(\frac{\Delta_t}{D}, 0, 1\right)\right) \right) & \text{otherwise, if it rained within } D \text{ hours},\\[4pt]
    \tau & \text{otherwise,}
  \end{cases}
\end{equation}

where $r_k$ is the rainfall at timestep $k$, $\theta = 0.1$ the small amount that
counts as wet, $m = 2.0$ the multiplier, $D = 36$ hours the post rain decay
window, and $\Delta_t$ the hours elapsed since the lookback window last closed
on wet weather. Thus, we double the bar while rain is recent and then ease it
back to normal over a day and a half rather than dropping it in one step. The
12 hour window and the doubling are the tuned values. They suppress the
routine first flush of a storm while leaving room to catch a deviation large
enough to clear twice the normal bar.

We added the taper because the step alone left a gap. With a hard edge, the
threshold returns to normal as soon as the lookback window clears, so a delayed
first flush pulse arriving a few hours later meets the ordinary bar and is
flagged as a discharge. The decay covers that period and returns the threshold
to normal smoothly instead.

The adjustment has a cost in the other direction. A discharge that occurs
during or shortly after rain has to clear a bar that is up to twice as high, so
real events in wet weather are harder to detect, and one confirmed event in the
catalog is missed for this reason. Setting $m$ and $D$ is a trade-off between
missing discharges in wet weather and flagging ordinary runoff, and it is a
judgment about what the system should flag rather than a flaw in the
mathematics.

\section{Deployment and Current Results}
\label{sec:deployment}

\subsection{Serving}
\label{sec:serving}

DuskCrayfish runs in the StrawberryWatch dashboard. It is the only model
switched on there by default, and it runs with two explainers attached,
Newtwork Run and Trial Bed, which are described in the support modules report.
Each live check reads the last two days of readings from every site and judges
the four nodes with the two rules of Section~\ref{sec:detection}.
Table~\ref{tab:serving} gives the calibrated quantities the served checkpoint
carries, all in normalized units.

\begin{table}[t]
  \centering
  \caption{Per-node calibration of the served checkpoint, in normalized units.
    The first three columns belong to the forecast residual rule
    (Equations~\ref{eq:score} and~\ref{eq:threshold}) and the last two to the
    level shift rule (Equation~\ref{eq:level-shift}).}
  \label{tab:serving}
  \begin{tabular}{lrrrrr}
    \toprule
    Node & Error median & Error IQR & $\tau_i$ & Level median & Level IQR \\
    \midrule
    north fork 0 & 0.0140 & 0.0169 & 12.091 & 0.540 & 0.252 \\
    south fork 2 & 0.0093 & 0.0140 & 11.817 & 1.433 & 0.661 \\
    south fork 1 & 0.0102 & 0.0139 & 12.756 & 0.000 & 1.424 \\
    oxford       & 0.0158 & 0.0250 & 2.445  & $-1.015$ & 0.204 \\
    \bottomrule
  \end{tabular}
\end{table}

Serving differs from training in three details. (1) Gaps of up to 3 hours are
filled in a live window, against the 24 hour limit in Table~\ref{tab:config}
that was used to build the training corpus. (2) The dashboard floors every
timestamp to its quarter hour before the model reads it, and averages readings
that land on the same quarter hour. (3) The rain adjustment of
Section~\ref{sec:rain} is expressed through the shared rain gate, with the same
12 hour lookback, 0.1 mm wet threshold and 36 hour taper, and the dashboard lets
an operator change the multiplier from its default of 2.0. A site can be
excluded from alerts but not from the model's input, since leaving a site out
of the input would require retraining without it.

The website also holds back a model's further alerts for a set time after it
sends one, 60 minutes by default and adjustable from 15 minutes to 24 hours.
We added this after the model sent repeated alerts during rain and cold
weather in October 2026. It limits how often an operator is alerted and does
not change what the model flags.

\subsection{Labelled events}
\label{sec:events}

On 25 October 2026 we replayed the four labelled anomalies in the event
catalog through the served detector, with its rain adjustment on.
Table~\ref{tab:events} gives the result at the labelled site's conductivity
node.

\begin{table}[t]
  \centering
  \caption{The labelled anomalies replayed through the served checkpoint.
    Steps over are counted against the rain-adjusted line, and a rule fires at
    3 steps. Peaks are given against the line before rain moved it.}
  \label{tab:events}
  \begin{tabular}{llrrrrl}
    \toprule
     & & \multicolumn{2}{c}{Forecast residual} & \multicolumn{2}{c}{Level shift} & \\
    Event & Site & Steps & Peak / $\tau_i$ & Steps & Peak / $k$ & Flagged \\
    \midrule
    June 2025 discharge      & south fork 1 & 3 & 27.031 / 12.756  & 0  & 1.190 / 4.0 & yes \\
    September 2025 overnight & south fork 1 & 8 & 116.342 / 12.756 & 0  & 1.075 / 4.0 & yes \\
    September 2025 overnight & south fork 2 & 1 & 29.148 / 11.817  & 0  & 3.541 / 4.0 & no \\
    March 2026 hydrant       & north fork 0 & 2 & 33.531 / 12.091  & 10 & 6.365 / 4.0 & yes \\
    \bottomrule
  \end{tabular}
\end{table}

The served checkpoint flagged three of the four. The June discharge crossed the
forecast line at exactly the three steps the rule needs. The hydrant break
crossed the forecast line at only two steps and was flagged by the level shift
rule, which held above its bar for ten steps, as Section~\ref{sec:rule2}
anticipates for a step that the forecast absorbs. The south fork 2 half of the
September event crossed the forecast line once, at a peak of more than twice
the line, and was not flagged because one step is below the three the rule
requires.

\section{Limitations and Future Work}
\label{sec:limitations}

Nearly every choice in the design points in the same direction: keep the model
small and matched to limited, rolling data, encode the known physics of water
flow directly in the graph rather than making the model learn it, and decide
anomalies on the one channel that carries the signal.

We have since addressed three of the limitations named in an earlier version
of this document. (1) We added cyclical time features, so the model now has
hour of day and day of year context and no longer reads an out of season window
as unusual everywhere. (2) The loss scores conductivity, depth and temperature
rather than conductivity alone. (3) We widened the post rain window into the
decay described in Section~\ref{sec:rain}. These were listed as natural next
steps, and we took them.

The remaining limitations form a shorter list, and the first two are the most
important.

Because detection scores conductivity alone, an anomaly that appears only in
another channel (e.g., a dissolved oxygen crash from an oil sheen) would not be
caught by the current alarm. The loss now watches three channels, but detection
still watches one.

The rain adjustment gives up some detection in order to suppress false alarms.
A discharge during or shortly after rain has to clear up to twice the ordinary
bar, and at least one confirmed event in the catalog is not detected under the
deployed multiplier. Lowering the multiplier would recover that event at the
cost of flagging ordinary runoff. We chose these values by judgment rather than
deriving them, and they belong to whoever owns the alerting budget.

The three-step rule that keeps one noisy reading from raising an alarm also
hides a brief real excursion. The south fork 2 half of the September 2025
event (Section~\ref{sec:events}) crossed the forecast line once, by more than
a factor of two, and was not flagged. Whether a single step that far over
should count is a judgment for whoever owns the alerting budget.

The network has four nodes, not five. We dropped footbridge for the coverage
reason given in Section~\ref{sec:footbridge}, which keeps training and
deployment consistent but leaves the north branch with no sensor between north
fork 0 and oxford. An event entering there is seen only once it reaches oxford.

The threshold rests on a single chronological split of a corpus that begins in
June 2025. The held-out portion covers a season the training portion does not
contain, so distribution shift and overfitting cannot be fully separated with
the data as it stands. We leave a same season holdout, which would separate
them, as an open question for future work.

Finally, we retrain the model on a rolling window rather than training it once
and freezing it, so every number in this document describes the current
checkpoint rather than the architecture in general. The thresholds in
particular are fitted quantities and will move when the model is refitted.

\bibliographystyle{unsrtnat}
\bibliography{references}

\end{document}
